\documentclass[11pt]{article}
\usepackage[margin=1in]{geometry}
\usepackage{graphicx}
\usepackage{booktabs}
\usepackage{adjustbox}
\usepackage{longtable}
\usepackage{amssymb}
\usepackage{mathtools}
\usepackage{float}

\newcommand{\SuppTableInput}[1]{%
  \begingroup
  \footnotesize
  \setlength{\tabcolsep}{4pt}%
  \renewcommand{\arraystretch}{1.05}%
  \noindent\centering
  \begin{adjustbox}{max width=\linewidth}
    \input{#1}%
  \end{adjustbox}
  \par
  \endgroup
}

\title{Supplementary Information:\\ Stable and interchangeable deformation pathways preserve pelvic mechanical function across anatomical variability}
\author{Petr Heny\v{s}, Michal Kucha\v{r}, Alexander Mor\'{a}vek, Niels Hammer}
\date{}

\begin{document}
\maketitle

\noindent
This Supplementary Information collects robustness, specificity, and model-boundary analyses that support the main manuscript.
Throughout, the childbirth-related load cases remain linear proxy loads rather than nonlinear labour simulations.

\section*{Supplementary Note 1 | Spectral cluster architecture}

The main manuscript focuses on the robust low-rank backbone and on a small number of recurrent clustered subspaces.
The tables below provide the complete bookkeeping needed to track those clustered regimes across the cohort, including reference-mode occupancy, swap frequencies, and morphology-axis diagnostics.

\section*{Table S1 | Detected near-degenerate rank clusters}
\SuppTableInput{../../analysis_outputs/tables/clusters.tex}

\section*{Table S2 | Rank-cluster to reference-mode mapping}
\SuppTableInput{../../analysis_outputs/tables/rank_cluster_ref_mapping.tex}

\section*{Table S3 | Per-mode swap rates}
\SuppTableInput{../../analysis_outputs/tables/swap_rates.tex}

\section*{Table S4 | Adjacent rank-pair swap diagnostics}
\SuppTableInput{../report/swap_pair_diagnostics.tex}

\section*{Table S5 | Cluster-level exchange classification}
\SuppTableInput{../report/veering_vs_mixing_classification.tex}

\section*{Table S6 | Morphology-axis permutation null test}
\SuppTableInput{../../analysis_outputs/tables/veering_null_test.tex}

\clearpage
\section*{Supplementary Note 2 | Cohort robustness and sex effects}

The main text uses sex effects only where they materially sharpen interpretation of the recurrent clustered subspaces, especially the 9--10 inlet-related pair.
The figure and tables below provide the broader sex-stratified summary and the bootstrap evidence that the dominant clusters are stable features of the cohort rather than products of a few outlying subjects.

\section*{Table S7 | Bootstrap robustness under random 80\% subcohorts}
\SuppTableInput{../../analysis_outputs/tables/cohort_bootstrap_robustness.tex}

\section*{Figure S1 | Sex-stratified routing distributions across all fifteen modes}
\begin{figure}[H]
  \centering
  \includegraphics[width=\linewidth]{../../analysis_outputs/plos_revision/figures/routing_distributions.pdf}
\end{figure}

\section*{Table S8 | Sex-stratified tests for spectral and subspace metrics}
\SuppTableInput{../../analysis_outputs/tables/sex_stratified_summary.tex}

\clearpage
\section*{Supplementary Note 3 | Load-routing specificity}

The main manuscript uses standing and proxy childbirth loads to show the contrast between a robust low-rank backbone and selectively recruited higher-rank subspaces.
The materials below provide the fuller load-specificity diagnostics, including cluster-level function-locking, dominant-mode sets, and pairwise overlap of full routing profiles.

\section*{Figure S2 | Load-function coupling correlations across prevalent clusters}
\begin{figure}[H]
  \centering
  \includegraphics[width=\linewidth]{../../analysis_outputs/plos_revision/figures/coupling_routing_heatmap.pdf}
\end{figure}

\section*{Table S9 | Load-function coupling synthesis}
\SuppTableInput{../../analysis_outputs/tables/function_locking_synthesis.tex}

\section*{Table S10 | Dominant mode sets for standing and proxy loading}
\SuppTableInput{../../analysis_outputs/tables/mode_activation_dominant_sets.tex}

\section*{Table S11 | Pairwise overlap of load-routing profiles}
\SuppTableInput{../../analysis_outputs/tables/mode_activation_overlap.tex}

\clearpage
\section*{Supplementary Note 4 | Model-boundary, sensitivity, and deformation atlas}

These analyses bound the main claims rather than extend them.
Allometry quantifies what is lost when the main spectral analysis is performed on size-normalized pelves.
Material uncertainty quantifies how much absolute eigenvalue identifiability would degrade under broader soft-tissue and density-law scatter, while showing that the relative fragile-versus-robust gap architecture is preserved.
Robustness checks confirm stability across percentile thresholds and reference subject choices.
Finally, the full 15-mode displacement field atlas is rendered on the common reference domain.

\section*{Figure S3 | Robustness and reconstruction error}
\begin{figure}[H]
  \centering
  \includegraphics[width=\linewidth]{../../analysis_outputs/plos_revision/figures/robustness.pdf}
\end{figure}

\section*{Figure S4 | Reference pelvic deformation atlas (fifteen modes)}
\begin{figure}[H]
  \centering
  \includegraphics[width=\linewidth,height=0.88\textheight,keepaspectratio]{../../analysis_outputs/plos_revision/figures/reference_modes.pdf}
\end{figure}

\clearpage
\section*{Figure S5 | Allometry dashboard}
\begin{figure}[H]
  \centering
  \includegraphics[width=\linewidth]{../../analysis_outputs/figures/allometry_dashboard.pdf}
\end{figure}

\section*{Table S12 | Allometry summary}
\input{../../analysis_outputs/tables/allometry_summary.tex}

\clearpage
\section*{Table S13 | Combined first-order material uncertainty summary}
\input{../../analysis_outputs/material_uncertainty/material_uncertainty_table.tex}

\end{document}
